\section{Series temporales económicas}

\subsection{Definición y relevancia}

Una serie temporal es una secuencia de observaciones ordenadas en el tiempo. En el software ICOCIV, cada observación corresponde al valor mensual de un índice para una categoría técnica seleccionada. Si \(y_t\) representa el índice en el mes \(t\), la serie se expresa como:

\begin{equation}
    \{y_1, y_2, \ldots, y_T\},
    \label{eq:serie_temporal}
\end{equation}

donde \(T\) es el número de periodos observados. La posición temporal de cada dato no es intercambiable. Cambiar el orden de los meses destruiría la información dinámica de la serie y haría inválido cualquier análisis predictivo.

En índices de construcción, el orden temporal contiene información económica. Un aumento sostenido en materiales puede reflejar inflación de insumos, restricciones logísticas o encarecimiento de materias primas. Una caída puntual puede responder a normalización de precios, cambios de oferta o ajustes metodológicos. Por tanto, el análisis estadístico debe preservar la secuencia cronológica, como recomienda la literatura clásica de forecasting para datos observados secuencialmente \parencite{hyndman2021fpp3}.

\subsection{Componentes de una serie económica}

Una serie mensual de índices puede descomponerse conceptualmente en varios componentes:

\begin{equation}
    y_t = T_t + S_t + C_t + \varepsilon_t,
    \label{eq:componentes_serie}
\end{equation}

donde \(T_t\) representa tendencia, \(S_t\) representa estacionalidad, \(C_t\) representa ciclos o fluctuaciones de mediano plazo y \(\varepsilon_t\) representa ruido o variación no explicada. Esta descomposición es conceptual: el software no implementa todavía un modelo formal de descomposición estacional, pero sí evalúa rasgos relacionados con tendencia, volatilidad y estabilidad.

La tendencia es especialmente relevante en índices de costos porque los precios de construcción tienden a incorporar procesos acumulativos. La estacionalidad puede aparecer si ciertos insumos o actividades presentan patrones recurrentes a lo largo del año. La volatilidad refleja la magnitud de los cambios mensuales, y puede aumentar en periodos de alta incertidumbre económica.

\subsection{Persistencia temporal}

La persistencia temporal se refiere a la tendencia de una serie a conservar parte de su comportamiento pasado. En precios e índices, un aumento no siempre se corrige rápidamente; puede transmitirse a meses posteriores por contratos, inventarios, expectativas o rigideces de mercado. Esta persistencia explica por qué la autocorrelación es frecuente en series económicas \parencite{hyndman2021fpp3}.

En términos simples, si \(y_t\) depende parcialmente de \(y_{t-1}\), entonces:

\begin{equation}
    y_t = \alpha + \rho y_{t-1} + u_t,
    \label{eq:persistencia}
\end{equation}

donde \(\rho\) mide el grado de persistencia y \(u_t\) es una perturbación. El software no usa esta ecuación como modelo principal de pronóstico, pero el concepto justifica la necesidad de pruebas como Durbin-Watson sobre residuos y backtesting temporal.

\subsection{Diferencia frente a datos independientes}

En una muestra transversal independiente, el orden de las observaciones suele no modificar el análisis. Por ejemplo, en una encuesta de viviendas, reorganizar las filas no cambia el significado estadístico básico. En una serie temporal, el orden es parte de los datos. Entrenar un modelo con información posterior para predecir un periodo anterior introduce fuga de información y produce una evaluación artificialmente optimista.

Esta diferencia afecta la validación predictiva. Dividir aleatoriamente una serie mensual de ICOCIV en entrenamiento y prueba permitiría que el modelo aprenda de meses futuros para predecir meses pasados. Eso contradice el escenario real de uso: cuando el usuario proyecta junio de 2026, solo dispone de información hasta el último mes publicado por el DANE. Por esta razón, el software adopta validación tipo \textit{walk-forward}, donde cada predicción usa únicamente datos anteriores.

\subsection{Estabilidad y cambios estructurales}

La estabilidad de una serie se refiere a que sus patrones estadísticos no cambien abruptamente. En índices de costos, la estabilidad puede romperse por choques económicos, cambios en metodología, crisis logísticas, variaciones fuertes de tasa de cambio o modificaciones en precios internacionales de insumos. Cuando ocurre un cambio estructural, un modelo ajustado al pasado puede perder capacidad predictiva.

El software no asume estabilidad perfecta. Por eso, sus reportes deben leerse como análisis condicionados a la información histórica disponible. La validación de calidad de datos y el backtesting ayudan a detectar señales de inestabilidad, pero no eliminan la incertidumbre inherente al forecasting económico.
